%========================== ORGANISME D'ACCUEIL =============================
\chapter{Présentation de l'organisme d'accueil}

\section{SBS Data Factory}

\textbf{SBS Data Factory} est l'entreprise au sein de laquelle s'est déroulé
ce stage de projet de fin d'année. Son activité s'inscrit dans le domaine de
la \textbf{donnée} et des \textbf{solutions numériques} appliquées au
marketing et à la publicité digitale. Dans ce cadre, l'entreprise s'intéresse
notamment aux dispositifs permettant de relier la présence digitale des
annonceurs à des résultats concrets et mesurables --- en particulier la
génération de \textbf{visites en magasin} (\emph{drive-to-store}).

Par souci de rigueur, ce rapport ne détaille pas d'informations
organisationnelles ou confidentielles internes à l'entreprise. Il se
concentre sur le \textbf{travail réalisé} durant le stage, à savoir la
conception et le développement d'un prototype de plateforme web.

\section{Cadre et sujet du stage}

Le stage a porté sur la réalisation d'un \textbf{prototype web de gestion de
campagnes drive-to-store}. L'objectif était de matérialiser, dans une
application fonctionnelle, l'ensemble du parcours d'un utilisateur de type
DSP : depuis la connexion jusqu'au reporting, en passant par la création de
campagne, le ciblage des magasins et la génération de créatives.

Le travail a été mené de manière \textbf{autonome et encadrée}, sous la
supervision de \textbf{Madame Dina TAALA}, avec un rythme de développement
\textbf{hebdomadaire} : chaque semaine correspondait à un ou plusieurs modules,
suivis d'une phase de validation et de documentation. Cette organisation, à la
fois progressive et traçable, est détaillée au chapitre consacré à la
méthodologie.

\section{Positionnement du prototype}

Il convient de préciser d'emblée le positionnement du livrable :

\begin{itemize}
  \item il s'agit d'un \textbf{prototype}, destiné à valider les parcours et
        l'architecture, et non d'un produit déployé en production ;
  \item les données sont majoritairement \textbf{simulées} (\emph{mockées} ou
        en mémoire), ce qui permet de faire fonctionner l'application sans
        infrastructure de base de données ;
  \item l'accent a été mis sur la \textbf{qualité logicielle} (structure du
        code, séparation des responsabilités, documentation) afin de faciliter
        une évolution ultérieure vers une version connectée à des données
        réelles.
\end{itemize}

Ce positionnement, transparent et assumé, est rappelé tout au long du rapport
et fait l'objet d'un chapitre dédié aux limites et perspectives.
